\documentclass{situationist-preprint}

\title{Retroactive Significance Without Retrocausation}
\msid{REL-22-a}
\msfamily{REL-22}
\msdiscipline{Philosophy of Religion}
\mstype{Principal manuscript}
\mskeywords{omnipotence; fixity of the past; significance; narrative sentences; defeat of evil; branching time; redemption}

\begin{document}
\maketitle

\begin{abstract}
Can God change the past? The tradition answers no, and we agree: to change the past is to bring it about that something that happened did not happen, which is impossible, and omnipotence extends only to what is possible. Nor does a timeless God affect the past by acting from a later time. But the tradition also speaks of redemption, of past suffering made good and past failures made the beginning of something. We argue that this language is literally true of a different object. Events have significance properties, which they have in virtue of their relations to later events, and what significance an event has is settled only as later events become settled. God, and creatures, can therefore change the settled significance of past events by what is later brought about, without changing anything that happened. We give the conditions under which such change is possible, distinguish it from retrocausation, and show that it allows past evils to be defeated, in Chisholm's sense, without ceasing to be evils.
\end{abstract}

\section{The problem}

Aristotle quotes Agathon as saying that this alone is denied even to God, to make undone what has been done \citep[VI.2, 1139b]{AristotleNE}. Aquinas agreed that God cannot make the past not to have been, since that would imply a contradiction, and contradictories fall outside the scope of omnipotence \citep[I, q.~25, a.~4]{AquinasST}. Peter Damian is often read as having held otherwise, but the dominant view has been that the past is beyond the reach of any power.

The religious traditions nonetheless speak as though something about the past can be altered. Suffering is said to be redeemed, a betrayal to become the occasion of reconciliation, a death to be given meaning by what follows. \citet{Adams1999} argued that horrendous evils can be defeated within a participant's life by their integration into a relationship with God, so that the participant would not, in retrospect, wish them away. \citet{Velleman1991} argued, independently of theology, that the value of a life is not the sum of the momentary well-being in it, and that a misfortune can be redeemed by what it leads to, so that a life's narrative shape matters to its value. Neither claims that redemption changes what happened.

We argue that the two strands are consistent. God cannot change the past, and does not causally affect it. But God can change the significance of past events in light of later ones, because significance is relational and is settled only as later events are.

\section{Framework}

\begin{definition}[Changing the past]\label{d:change}
To \emph{change the past} is to bring it about, at or from some time, that an event which occurred at an earlier time $t$ did not occur at $t$, or that a fact about $t$ which obtained does not obtain.
\end{definition}

\begin{definition}[Significance]\label{d:sig}
A \emph{significance property} of an event $e$ is a property $e$ has in virtue of its relations to later events: being the first battle of a long war, being the occasion of a reconciliation, being part of a life that went well on the whole.
\end{definition}

Significance properties are those ascribed by what \citet[ch.~8]{Danto1965} called narrative sentences, which describe an earlier event by reference to a later one, as in \enquote{the Thirty Years' War began in 1618}, a sentence no one in 1618 could have known to be true.

\begin{assumption}[Branching]\label{a:branch}
Times are ordered in a branching structure. A proposition about the future is settled at $t$ if it is true on every history through $t$. Future contingents are not settled, and lack a truth value at times before the matter is decided; tenseless propositions are true or false relative to the actual history.
\end{assumption}

Assumption~\ref{a:branch} is the view of ACT-17. It is compatible with the reality of the past defended in TSM-03 and with the B-theoretic account of time of TSM-01, applied to the actual history.

\begin{definition}[Settled significance]\label{d:settled}
Event $e$ has significance property $F$ \emph{as of} $t$ if $e$ has $F$ on every history through $t$.
\end{definition}

\section{The argument}

\begin{theorem}\label{t:nochange}
No agent, God included, can change the past.
\end{theorem}

\begin{proof}
Suppose an agent brought it about that an event $e$ that occurred at $t$ did not occur at $t$. Then $e$ both occurred and did not occur at $t$. No possible state of affairs answers to this description, and omnipotence, on the account of REL-04 and REL-02-a, is the ability to bring about every possible state of affairs that meets the stratification condition; an impossible state is no limitation on it. The same holds for any fact about $t$.
\end{proof}

Theorem~\ref{t:nochange} does not exclude affecting the past without changing it. TSM-04 argues that backward time travel is consistent provided the past is not changed: a traveller may cause events at an earlier time, which then were always as she caused them to be; \citet{Lewis1976} drew the distinction between changing the past and affecting it. God, however, is timeless according to REL-19, and acts from no time; God's willing what happens at $t$ is not later than $t$, and the timeless dependence of what happens on God's will is not backward causation. Neither form of influence on the past is a change in it, and neither is what the language of redemption claims.

\begin{lemma}\label{l:gain}
An event can lack a significance property as of $t_1$ and have it as of a later $t_2$, although nothing about the event, and nothing about any time up to $t_1$, differs between $t_1$ and $t_2$.
\end{lemma}

\begin{proof}
Let $e$ occur before $t_1$, and let $F$ be the property of being the occasion of a reconciliation that occurs at $t_2$. At $t_1$ the reconciliation is a future contingent, so by Assumption~\ref{a:branch} there are histories through $t_1$ on which it does not occur, and $e$ lacks $F$ as of $t_1$ by Definition~\ref{d:settled}. If the reconciliation occurs at $t_2$, every history through $t_2$ contains it, and $e$ has $F$ as of $t_2$. The histories through $t_2$ agree with those through $t_1$ on everything up to $t_1$, so nothing about $e$ or about any time up to $t_1$ differs.
\end{proof}

\begin{theorem}\label{t:sig}
God can change the significance of past events in light of later ones, without changing the past.
\end{theorem}

\begin{proof}
By Lemma~\ref{l:gain}, the settled significance of a past event changes when later events occur that stand in the relevant relations to it, and the change involves no difference in the event or in anything earlier. An omnipotent being can bring about, or cooperate in bringing about, later events of the relevant kinds, such as reconciliations, restorations and completions; so God can bring it about that a past event comes to have, as of a later time, a significance it lacked. By Theorem~\ref{t:nochange}, this is not a change in the past.
\end{proof}

Theorem~\ref{t:sig} does not require that God bring about the later events single-handedly. Many of the events that change the significance of the past are the free actions of creatures, an apology or a forgiveness, and on the view of REL-02-a some of the goods that matter most are those God does not settle. Creatures change the significance of their own pasts, and of one another's; PID-09 concludes that narrative identity is real and co-authored by others, which is a case of the same phenomenon.

\section{Defeat}

\citet{Chisholm1968} distinguished the \emph{balancing off} of an evil, in which a whole containing it contains goods that outweigh it, from its \emph{defeat}, in which the whole containing it is better than it would have been without it, because of the way the evil is related to the goods. The distinction depends on Moore's principle of organic unities, that the value of a whole need not be the sum of the values of its parts \citep[§18]{Moore1903}.

\begin{corollary}\label{c:defeat}
A past evil can come to be defeated, as of a later time, without ceasing to be an evil.
\end{corollary}

\begin{proof}
Being part of a whole whose value it enhances is a significance property in the sense of Definition~\ref{d:sig}, since the whole includes later events. By Theorem~\ref{t:sig}, an evil can come to have it as of a later time. Its intrinsic badness is a fact about it, not about its relations to later events, and by Theorem~\ref{t:nochange} that fact cannot be changed.
\end{proof}

Corollary~\ref{c:defeat} gives a precise sense to Adams's claim that horrendous evils can be defeated within a life. It also marks the limit of that claim: the suffering remains suffering, and the wrong done remains a wrong. What changes is what the suffering is part of. The companion manuscript REL-22-b examines how the language of redemption should be read in this light.

\section{Objections and replies}

\begin{objection}[Mere Cambridge change]
A change in an event's relations to later events is a Cambridge change, like Socrates's becoming shorter than Theaetetus as Theaetetus grows. It is not a real change in the event, and nothing of religious importance can rest on it.
\end{objection}

\begin{reply}
It is a Cambridge change in the event, and we claim no more. Whether it is important depends on whether significance properties matter, and they do: the value of a life, on Velleman's account, depends on its narrative relations, and the fitting response to a past event, regret or gratitude, depends on what it led to. That a property is relational does not make it trivial; being a parent is relational too.
\end{reply}

\begin{objection}[Tenseless truth]
On the B-theory, \enquote{$e$ is the occasion of a reconciliation} is tenselessly true or false. If it is true, it was always true, and nothing changes.
\end{objection}

\begin{reply}
The tenseless proposition is true relative to the actual history, and does not change its truth value. What changes is whether the significance is settled, by Definition~\ref{d:settled}, and under Assumption~\ref{a:branch} that is a real difference between times: before the reconciliation, no fact made it true at that time that $e$ would occasion one. If the future were settled, significance would not change; it would merely be discovered. The thesis depends on the openness of the future, which ACT-17 defends.
\end{reply}

\begin{objection}[Justification of evil]
If later goods can redeem earlier evils, the evils are justified by them, and whoever caused them is excused.
\end{objection}

\begin{reply}
Defeat is not justification. That an evil comes to be part of a whole it enhances does not show that anyone was permitted to cause it, and the agent who caused it remains responsible for what she did, since what she did, by Theorem~\ref{t:nochange}, is unchanged.
\end{reply}

\section{Conclusion}

God cannot change the past, since that would be to make what happened not have happened, and God does not reach the past by acting from a later time. God can change the significance of past events, because significance consists in relations to later events and is settled only as those events become settled. Redemption, so understood, alters what the past is part of, not what the past was.

\begin{thebibliography}{99}\small
\bibitem[Adams(1999)]{Adams1999} Adams, M.~M. (1999). \emph{Horrendous Evils and the Goodness of God}. Cornell University Press.
\bibitem[Aquinas(ST)]{AquinasST} Aquinas, T. \emph{Summa Theologiae}. Cited by part, question and article.
\bibitem[Aristotle(NE)]{AristotleNE} Aristotle. \emph{Nicomachean Ethics}. Cited by book, chapter and Bekker page.
\bibitem[Chisholm(1968)]{Chisholm1968} Chisholm, R.~M. (1968). The defeat of good and evil. \emph{Proceedings and Addresses of the American Philosophical Association}, 42, 21--38.
\bibitem[Danto(1965)]{Danto1965} Danto, A.~C. (1965). \emph{Analytical Philosophy of History}. Cambridge University Press.
\bibitem[Lewis(1976)]{Lewis1976} Lewis, D. (1976). The paradoxes of time travel. \emph{American Philosophical Quarterly}, 13(2), 145--152.
\bibitem[Moore(1903)]{Moore1903} Moore, G.~E. (1903). \emph{Principia Ethica}. Cambridge University Press.
\bibitem[Velleman(1991)]{Velleman1991} Velleman, J.~D. (1991). Well-being and time. \emph{Pacific Philosophical Quarterly}, 72(1), 48--77.
\end{thebibliography}

\end{document}
